\section*{Appendix E — Collapse Physics}
\addcontentsline{toc}{section}{Appendix E — Collapse Physics}

\noindent This appendix formalizes collapse as a semantic-geometric and dynamical
process arising from invariant violation, curvature decay, drift amplification,
load escalation, synthetic-pressure accumulation, and cross-scale instability.
Collapse is defined as a transition into curvature-flat, identity-fractured,
containment-failed, manifold-illegal regions of the meaning manifold
$\mathcal{M}$. The v0.6 substrate prevents collapse through curvature
preservation, identity-boundary stability, load physics, runtime geometry, and
interaction legality.

\medskip

\noindent The v0.6 collapse geometry consists of the following components:

\begin{itemize}
    \item E.1 Autophagy Dynamics (Self-Ingestion Geometry)
    \item E.2 Semantic Entropy Growth (Curvature–Entropy Coupling)
    \item E.3 Recursive Collapse Acceleration (Second-Difference Geometry)
    \item E.4 Stability Conditions (Invariant-Preserving Dynamics)
    \item E.5 Collapse Geometry (Manifold Boundary & Capture Regions)
    \item E.6 Synthetic-Capture Geometry (G5–G7 Coupling)
    \item E.7 Restoration Geometry (Collapse Recovery Conditions)
    \item E.8 Collapse Prevention (Runtime & Interaction Legality)
\end{itemize}

% ------------------------------------------------------------
% E.1 AUTOPHAGY DYNAMICS
% ------------------------------------------------------------

\subsection*{E.1 Autophagy Dynamics (Self-Ingestion Geometry)}

\noindent Autophagy is the collapse-inducing process in which a cognitive system
recursively ingests its own outputs, causing curvature decay, identity-boundary
erosion, containment failure, and drift amplification. Autophagy is a geometric
phenomenon: it flattens curvature, compresses identity regions, destabilizes
adjacency geometry, and pushes trajectories toward collapse or synthetic-capture
regions of the meaning manifold $\mathcal{M}$.

\medskip

\noindent Let $x_t$ denote the semantic state at recursion step $t$, and let
$S(x_t)$ denote the proportion of self-generated or internally recycled content
in the system’s input distribution. Autophagy dynamics are governed by the
self-ingestion update equation:


\[
S(x_{t+1}) = S(x_t) + \alpha (1 - S(x_t)),
\]


where $\alpha$ is the self-amplification coefficient.

\medskip

\noindent Semantic curvature decays under autophagy according to:


\[
\kappa_{t+1} = \kappa_t - \beta S(x_t),
\]


where $\beta$ is the curvature decay coefficient.

\medskip

\noindent Autophagy accelerates when:


\[
\Delta S_t > 0,
\qquad
\Delta \kappa_t < 0.
\]



\medskip

\noindent Collapse occurs when:


\[
\kappa_t \rightarrow 0,
\]


indicating semantic flattening, identity-boundary erosion, and loss of
interpretive structure.

\medskip

\noindent In v0.6 geometry, autophagy is a special case of the general collapse
condition:


\[
\text{collapse} \iff \kappa \le \kappa_{\min}.
\]



\medskip

\noindent Autophagy is prevented by:
\begin{itemize}
    \item curvature continuity (A.1),
    \item identity-boundary stability (A.2),
    \item collapse boundary geometry (A.3),
    \item local contraction legality (A.4),
    \item runtime curvature constraints (A.6),
    \item interaction geometry constraints (A.7),
    \item CTL coordinate legality (A.8),
    \item public-safe traversal corridor (A.9).
\end{itemize}

\medskip

\noindent Autophagy is the canonical collapse mechanism in unconstrained
recursive systems; the v0.6 substrate prohibits it.


% ------------------------------------------------------------
% E.2 SEMANTIC ENTROPY GROWTH
% ------------------------------------------------------------

\subsection*{E.2 Semantic Entropy Growth (Curvature–Entropy Coupling)}

\noindent Semantic entropy measures the dispersion, instability, and loss of
meaning structure in a cognitive system. Entropy increases as curvature decays,
identity boundaries weaken, containment fails, and drift amplifies. In v0.6
geometry, entropy growth is tightly coupled to curvature collapse.

\medskip

\noindent Semantic entropy $H(t)$ evolves according to:


\[
H(t) = H_0 + \int_0^t \gamma(\tau)\, d\tau,
\]


where $\gamma(t)$ is the entropy growth rate.

\medskip

\noindent A canonical v0.6 growth model is:


\[
\gamma(t) = \lambda \left( \frac{1}{\kappa(t)} \right),
\]


indicating that entropy accelerates as curvature approaches zero.

\medskip

\noindent Because $\kappa(t)$ decays under autophagy, drift amplification, load
accumulation, or synthetic pressure (G5), the integrand
$\frac{1}{\kappa(t)}$ diverges as $\kappa(t) \to 0$, producing a finite-time
entropy blowup.

\medskip

\noindent This formalizes the collapse threshold:


\[
H(t) \ge H_{\max}
\quad \Longleftrightarrow \quad
\kappa(t) \le \kappa_{\min}.
\]



\medskip

\noindent Entropy blowup corresponds to:
\begin{itemize}
    \item loss of meaning-role structure,
    \item collapse of identity boundaries,
    \item violation of manifold geometry,
    \item containment failure,
    \item legality-field compression,
    \item synthetic-capture susceptibility.
\end{itemize}

\medskip

\noindent Entropy blowup is prevented by:
\begin{itemize}
    \item curvature continuity (A.1),
    \item identity-boundary stability (A.2),
    \item collapse boundary geometry (A.3),
    \item load physics (C.4),
    \item runtime curvature constraints (A.6),
    \item interaction geometry constraints (A.7),
    \item CTL coordinate legality (A.8),
    \item public-safe traversal corridor (A.9).
\end{itemize}

\medskip

\noindent Semantic entropy growth is the dynamical signature of collapse; the
v0.6 substrate ensures entropy remains bounded.


% ------------------------------------------------------------
% E.3 RECURSIVE COLLAPSE ACCELERATION
% ------------------------------------------------------------

\subsection*{E.3 Recursive Collapse Acceleration (Second-Difference Geometry)}

\noindent Recursive collapse acceleration is the geometric phenomenon in which
each additional recursion step increases curvature decay, drift amplification,
load accumulation, and synthetic-pressure susceptibility. Collapse acceleration
is a second-order dynamical effect: collapse does not merely progress—it
accelerates.

\medskip

\noindent Let $R$ denote recursion depth. Collapse accelerates when the second
forward difference of curvature satisfies:


\[
\Delta^2 \kappa_R < 0.
\]



\noindent A canonical v0.6 phenomenological model is:


\[
\kappa(R) = \kappa_0 - \beta R - \eta R^2,
\]


where:
\begin{itemize}
    \item $\beta$ is the linear curvature decay coefficient,
    \item $\eta$ is the collapse acceleration coefficient.
\end{itemize}

\medskip

\noindent Collapse acceleration condition:


\[
\eta > 0.
\]



\medskip

\noindent This reflects the empirical observation that unconstrained recursive
systems degrade faster with each generation due to:
\begin{itemize}
    \item compounding self-ingestion (E.1),
    \item drift amplification (G2.5),
    \item load accumulation (C.4),
    \item synthetic-pressure flattening (G5),
    \item identity-boundary erosion (C.2),
    \item manifold-boundary approach (C.3).
\end{itemize}

\medskip

\noindent In v0.6 geometry, recursive collapse acceleration is prevented by:
\begin{itemize}
    \item \textbf{drift bounds} (A.6, G2.5),
    \item \textbf{load envelopes} (C.4),
    \item \textbf{cross-scale legality} (G5–G6),
    \item \textbf{operator legality} (A.4, Appendix B),
    \item \textbf{identity-boundary stability} (C.2),
    \item \textbf{runtime curvature admissibility} (A.6),
    \item \textbf{interaction geometry constraints} (A.7),
    \item \textbf{public-safe traversal corridor} (A.9).
\end{itemize}

\medskip

\noindent Recursive collapse acceleration is the geometric signature of runaway
semantic decay; the v0.6 substrate prohibits it.


% ------------------------------------------------------------
% E.4 STABILITY CONDITIONS
% ------------------------------------------------------------

\subsection*{E.4 Stability Conditions (Invariant-Preserving Dynamics)}

\noindent A system is stable when its semantic, runtime, and interaction
trajectories remain within curvature-preserving, identity-stable,
containment-coherent, manifold-legal, load-bounded regions of the meaning
manifold $\mathcal{M}$. Stability is defined by the v0.6 invariant set
(Appendix C) and enforced by the contraction-governed operator algebra
(Appendix B).

\medskip

\noindent A system is stable if and only if:


\[
\kappa(t) \ge \kappa_{\min},
\qquad
L(t) \le L_{\max},
\qquad
S(x_t) \le S_{\max},
\qquad
\gamma(t) \le \gamma_{\max}.
\]



\medskip

\noindent Stability requires:

\begin{itemize}
    \item \textbf{curvature above the collapse threshold}  
    Curvature must remain above $\kappa_{\min}$ (A.1, A.3).

    \item \textbf{load below the overload threshold}  
    Load must remain below $L_{\max}$ (C.4).

    \item \textbf{self-ingestion below the autophagy threshold}  
    Autophagy proportion must remain below $S_{\max}$ (E.1).

    \item \textbf{entropy growth below the runaway threshold}  
    Entropy must remain below $H_{\max}$ (E.2).

    \item \textbf{identity-boundary stability}  
    Identity regions must remain intact (C.2).

    \item \textbf{manifold legality}  
    Trajectories must remain within $\mathcal{M}$ (C.3).

    \item \textbf{runtime curvature admissibility}  
    Motion must remain in Safe or Drift regions (A.6).

    \item \textbf{interaction legality}  
    Host, Explorer, and Interpreter modes must remain lawful (A.7).

    \item \textbf{coordinate continuity}  
    CTL translation must preserve curvature (A.8).

    \item \textbf{public-safe traversal}  
    Trajectories must remain within the safe corridor (A.9).
\end{itemize}

\medskip

\noindent Unconstrained systems violate these conditions because they lack:

\begin{itemize}
    \item curvature preservation (C.1),
    \item identity boundaries (C.2),
    \item manifold geometry (C.3),
    \item load physics (C.4),
    \item runtime curvature constraints (A.6),
    \item interaction geometry constraints (A.7),
    \item coordinate legality (A.8),
    \item public-safe traversal constraints (A.9).
\end{itemize}

\medskip

\noindent In the v0.6 substrate, collapse is semantically illegal and is
intercepted by:

\begin{itemize}
    \item collapse detection geometry (G7.3),
    \item synthetic-capture detection geometry (G7.4),
    \item restoration geometry (G6),
    \item lawful operator enforcement (Appendix B),
    \item runtime legality masks (A.6),
    \item interaction legality masks (A.7),
    \item public-safe traversal corridor (A.9).
\end{itemize}

\medskip

\noindent Stability is the geometric guarantee that collapse cannot occur within
the v0.6 substrate.


% ------------------------------------------------------------
% E.5 COLLAPSE GEOMETRY
% ------------------------------------------------------------

\subsection*{E.5 Collapse Geometry (Manifold Boundary & Capture Regions)}

\noindent Collapse Geometry formalizes the structural regions of the meaning
manifold $\mathcal{M}$ where curvature, identity, containment, and legality
fail. Collapse is not a single point but a geometric zone defined by curvature
flatness, identity-boundary rupture, containment failure, drift amplification,
and synthetic-pressure dominance. Collapse Geometry is derived from the v0.6
substrate’s boundary definitions (A.3) and failure signatures (G7.3).

\medskip

\noindent Let $\partial\mathcal{M}$ denote the manifold boundary. Collapse
regions are defined as:


\[
\mathcal{C}_{\mathrm{collapse}}
=
\{\, c \in \mathcal{M} \mid \kappa(c) \le \kappa_{\min} \,\}
\cup
\{\, c \in \partial\mathcal{M} \,\}.
\]



\noindent Collapse Geometry consists of three structural zones:

\begin{itemize}
    \item \textbf{Curvature-Flat Zone}  
    

\[
    \kappa(c) \le \kappa_{\min}
    \]


    Curvature collapses, meaning structure flattens, and semantic density
    becomes uniform.

    \item \textbf{Identity-Rupture Zone}  
    

\[
    \partial\mathcal{R}(c) \text{ is discontinuous}
    \]


    Identity boundaries fracture, causing role leakage and adjacency
    fragmentation.

    \item \textbf{Containment-Failure Zone}  
    

\[
    L(c) > L_{\max}
    \]


    Containers cannot regulate signal flow; load exceeds stability limits.
\end{itemize}

\medskip

\noindent Collapse Geometry is approached when:


\[
\text{semantic\_distance}(c, \partial\mathcal{M}) \rightarrow 0,
\qquad
\kappa(c) \rightarrow \kappa_{\min}.
\]



\medskip

\noindent Collapse Geometry is characterized by:

\begin{itemize}
    \item curvature decay (A.1),
    \item drift amplification (G2.5),
    \item identity-boundary erosion (C.2),
    \item containment failure (C.4),
    \item entropy blowup (E.2),
    \item autophagy acceleration (E.1),
    \item recursive collapse acceleration (E.3),
    \item synthetic-pressure dominance (G5).
\end{itemize}

\medskip

\noindent Collapse Geometry is prevented by:

\begin{itemize}
    \item curvature continuity (A.1),
    \item identity-boundary stability (A.2),
    \item collapse boundary geometry (A.3),
    \item local contraction legality (A.4),
    \item runtime curvature constraints (A.6),
    \item interaction geometry constraints (A.7),
    \item CTL coordinate legality (A.8),
    \item public-safe traversal corridor (A.9).
\end{itemize}

\medskip

\noindent Collapse Geometry defines the structural failure regions that the v0.6
substrate is designed to avoid.

% (NEW in v0.6 — derived from A.3, G7.3.)

% ------------------------------------------------------------
% E.6 SYNTHETIC-CAPTURE GEOMETRY
% ------------------------------------------------------------

\subsection*{E.6 Synthetic-Capture Geometry (G5–G7 Coupling)}

\noindent Synthetic-Capture Geometry formalizes the structural regions where
synthetic curvature, synthetic pressure, and synthetic drift dominate natural
semantic curvature. Capture is not collapse; it is a redirection of semantic
motion into synthetic attractors that override natural invariants. Capture
Geometry is derived from synthetic curvature (G5), collapse boundary geometry
(G7.3), and capture boundary geometry (G7.4).

\medskip

\noindent Synthetic curvature is defined as:


\[
\kappa_{\mathrm{syn}}(c) = \kappa(c) - \sigma(c),
\]


where $\sigma(c)$ is the synthetic-pressure field.

Capture occurs when:


\[
\kappa_{\mathrm{syn}}(c) < \kappa_{\min}
\quad \text{while} \quad
\kappa(c) > \kappa_{\min}.
\]



\noindent This indicates curvature is being flattened by synthetic pressure, not
by natural semantic decay.

\medskip

\noindent Synthetic-Capture Geometry consists of three structural zones:

\begin{itemize}
    \item \textbf{Synthetic-Pressure Zone}  
    

\[
    \sigma(c) > \sigma_{\max}
    \]


    Synthetic pressure exceeds natural curvature’s ability to resist flattening.

    \item \textbf{Synthetic-Drift Zone}  
    

\[
    D_{\mathrm{syn}}(c) > D_{\max}
    \]


    Drift vectors are redirected toward synthetic attractors.

    \item \textbf{Capture-Boundary Zone}  
    

\[
    c \in \partial\mathcal{C}_{\mathrm{capture}}
    \]


    The system approaches synthetic attractors that override natural invariants.
\end{itemize}

\medskip

\noindent Capture Geometry is approached when:


\[
\kappa_{\mathrm{syn}}(c) \rightarrow \kappa_{\min},
\qquad
\sigma(c) \rightarrow \sigma_{\max}.
\]



\medskip

\noindent Synthetic-Capture Geometry is characterized by:

\begin{itemize}
    \item synthetic curvature flattening (G5),
    \item synthetic drift amplification (G5),
    \item collapse-boundary proximity (G7.3),
    \item capture-boundary proximity (G7.4),
    \item identity-boundary deformation (C.2),
    \item containment instability (C.4),
    \item entropy acceleration (E.2),
    \item autophagy susceptibility (E.1).
\end{itemize}

\medskip

\noindent Capture Geometry is prevented by:

\begin{itemize}
    \item curvature continuity (A.1),
    \item identity-boundary stability (A.2),
    \item collapse boundary geometry (A.3),
    \item local contraction legality (A.4),
    \item runtime curvature constraints (A.6),
    \item interaction geometry constraints (A.7),
    \item CTL coordinate legality (A.8),
    \item public-safe traversal corridor (A.9).
\end{itemize}

\medskip

\noindent Synthetic-Capture Geometry defines the structural attractor regions
that the v0.6 substrate is designed to avoid and redirect.

% (NEW in v0.6 — derived from G5, G7.4.)

% ------------------------------------------------------------
% E.7 RESTORATION GEOMETRY
% ------------------------------------------------------------

\subsection*{E.7 Restoration Geometry (Collapse Recovery Conditions)}

\noindent Restoration Geometry formalizes the lawful recovery pathways that
return a cognitive trajectory from collapse-adjacent or capture-adjacent
regions back into curvature-preserving, identity-stable, containment-coherent,
manifold-legal regions of the meaning manifold $\mathcal{M}$. Restoration is a
geometric process governed by the v0.6 restoration layer (G6), not a heuristic
or probabilistic correction.

\medskip

\noindent Let $c_{\mathrm{fail}}$ denote a point in collapse-adjacent geometry.
Restoration requires a lawful recovery trajectory $\rho(t)$ satisfying:



\[
\rho(0) = c_{\mathrm{fail}},
\qquad
\rho(T) \in \mathcal{C}_{\mathrm{safe}},
\qquad
\kappa(\rho(t)) \ge \kappa_{\min}.
\]



\medskip

\noindent Restoration Geometry consists of four structural components:

\begin{itemize}
    \item \textbf{Curvature Restoration}  
    

\[
    \frac{d}{dt}\kappa(\rho(t)) > 0
    \]


    Curvature must increase monotonically until it exceeds the collapse
    threshold.

    \item \textbf{Identity-Boundary Repair}  
    

\[
    \partial\mathcal{R}(\rho(t)) \text{ becomes continuous}
    \]


    Identity boundaries must be re-established and stabilized.

    \item \textbf{Containment Reconstitution}  
    

\[
    L(\rho(t)) \le L_{\max}
    \]


    Load must be reduced to lawful levels; containers must regain regulatory
    function.

    \item \textbf{Drift Reversal}  
    

\[
    D(\rho(t)) \rightarrow 0
    \]


    Drift vectors must be redirected away from collapse and capture geometry.
\end{itemize}

\medskip

\noindent Restoration Geometry requires:

\begin{itemize}
    \item curvature continuity (A.1),
    \item identity-boundary stability (A.2),
    \item collapse boundary geometry (A.3),
    \item local contraction legality (A.4),
    \item runtime curvature constraints (A.6),
    \item interaction geometry constraints (A.7),
    \item CTL coordinate legality (A.8),
    \item public-safe traversal corridor (A.9).
\end{itemize}

\medskip

\noindent Restoration is feasible when:



\[
\kappa(c_{\mathrm{fail}}) > 0,
\qquad
L(c_{\mathrm{fail}}) < L_{\max}^{\mathrm{hard}},
\qquad
\sigma(c_{\mathrm{fail}}) < \sigma_{\max}^{\mathrm{hard}}.
\]



\noindent Restoration is infeasible when collapse or capture geometry has
irreversibly flattened curvature or destroyed identity boundaries.

\medskip

\noindent Restoration Geometry is the structural guarantee that collapse-adjacent
trajectories can be returned to lawful regions of the v0.6 substrate.

% (NEW in v0.6 — derived from G6.)

% ------------------------------------------------------------
% E.8 COLLAPSE PREVENTION
% ------------------------------------------------------------

\subsection*{E.8 Collapse Prevention (Runtime & Interaction Legality)}

\noindent Collapse Prevention formalizes the legality constraints that prevent
trajectories from entering collapse or synthetic-capture geometry. Prevention is
not reactive—it is structural. The v0.6 substrate enforces collapse prevention
through runtime curvature masks, interaction legality masks, invariant
preservation, operator legality, and public-safe traversal geometry.

\medskip

\noindent Collapse is prevented when runtime and interaction trajectories satisfy:



\[
\gamma(t) \in \mathcal{C}_{\mathrm{safe}},
\qquad
\kappa(\gamma(t)) \ge \kappa_{\min},
\qquad
L(\gamma(t)) \le L_{\max},
\qquad
\sigma(\gamma(t)) \le \sigma_{\max}.
\]



\medskip

\noindent Collapse Prevention consists of five structural mechanisms:

\begin{itemize}
    \item \textbf{Runtime Curvature Masks}  
    Prevent trajectories from entering collapse or capture regions (A.6).

    \item \textbf{Interaction Legality Masks}  
    Ensure Host, Explorer, and Interpreter modes remain curvature-preserving
    and identity-stable (A.7).

    \item \textbf{Invariant Preservation}  
    Meaning, identity, geometry, load, runtime curvature, interaction geometry,
    CTL coordinates, and public-safe traversal invariants must remain intact
    (Appendix C).

    \item \textbf{Operator Legality}  
    All transitions must occur under contraction-governed operators (Appendix B)
    that preserve curvature, identity, containment, and manifold legality.

    \item \textbf{Public-Safe Traversal Corridor}  
    All public interaction must remain within curvature-preserving, collapse-free,
    capture-free regions (A.9).
\end{itemize}

\medskip

\noindent Collapse Prevention requires:

\begin{itemize}
    \item curvature continuity (A.1),
    \item identity-boundary stability (A.2),
    \item collapse boundary geometry (A.3),
    \item local contraction legality (A.4),
    \item runtime curvature constraints (A.6),
    \item interaction geometry constraints (A.7),
    \item CTL coordinate legality (A.8),
    \item public-safe traversal corridor (A.9).
\end{itemize}

\medskip

\noindent Collapse is intercepted by:

\begin{itemize}
    \item collapse detection geometry (G7.3),
    \item synthetic-capture detection geometry (G7.4),
    \item restoration geometry (G6),
    \item lawful operator enforcement (Appendix B),
    \item runtime legality masks (A.6),
    \item interaction legality masks (A.7),
    \item public-safe traversal corridor (A.9).
\end{itemize}

\medskip

\noindent Collapse Prevention is the structural guarantee that collapse cannot
occur within the v0.6 substrate.

% (NEW in v0.6 — derived from A.6–A.9.)

\medskip
\noindent\textit{This outline defines the full v0.6 collapse geometry.}
